%!TEX root = ../../thesis.tex

\section{An EfficientZero Implementation in Rust}\label{sec:impl-efficientzero}

In addition to the fuzzer, we provide an EfficientZero implementation in Rust
using the Burn deep learning framework.
The implementation provides a flexible interface which enables training agents for
various environments as well as a unified way to specify parameters and a dashboard
to monitor training progress.

\subsection{The Environment Model}

EfficientZero is an RL algorithm and therefore formally interacts with an MDP.
Making use of the rich type system of Rust, we model the MDP via an \code{Environment} trait
with associated \code{Action} and \code{Observation} types.
Actions and Observations must be representable as a list of floats for model inference and training.
The environment supports stepping its current state using an action and then returns
a reward and observation to the caller, as well as a flag \code{done} that indicates
that the episode has finished.
It also supports resetting its state to the beginning which returns an initial observation.

\begin{rust}[caption={The environment interface of the EfficientZero
    implementation.},label={lst:ez-abstraction}]
pub trait Observation: Clone {
    fn to_vec(&self) -> Vec<f32>;
    fn size() -> usize;
}

pub trait Action: Into<usize> + From<usize> + Copy {
    fn size() -> usize;
}

pub struct StepResult<O> {
    pub reward: f32,
    pub done: bool,
    pub observation: O,
}

pub trait Environment {
    type Action: Action;
    type Observation: Observation;

    fn step(&mut self, action: Self::Action) -> StepResult<Self::Observation>;
    fn reset(&mut self) -> Self::Observation;
    fn legal_actions() -> impl Iterator<Item = Self::Action>;
}
\end{rust}

Both associated types are constrained by nothing but what the networks need
from them.
An \code{Observation} must flatten into a vector of \code{f32} whose length is
fixed and known without an instance at hand, since it is the input width of the
representation network $h$.
An \code{Action} must map onto the index range $0, \dots, n - 1$, where $n$ is
the number of categories of the policy head $p$; the conversions to and from
\code{usize} are what allows the tree search and the replay buffer to store
actions as plain integers.
Note that \code{legal\_actions} is an associated function rather than a method
on the environment: beyond the root, \gls{mcts} descends through latent states
that have no environment behind them~(\cref{sec:ez-mcts}), so the set of
available actions cannot depend on the current state.

The environments for training and evaluating an agent are usually
different.
In the case of fuzzing, an evaluation environment might have different hooks
attached to extract information from the fuzzing campaign compared to a
training environment.
We therefore keep construction out of \code{Environment} itself and put it into
the two subtraits \code{TrainingEnvironment} and \code{EvaluationEnvironment},
each providing its own constructor.
Training is generic over the former, evaluation over the latter, and a target
that implements both can be trained and evaluated by the same driver.

\subsection{The Learner}

The \code{Learner} makes up the core of the implementation.
It contains an implementation of batched \gls{mcts}, the replay buffer, various
workers, and hook firing.

\subsubsection{MCTS}

\gls{mcts} was introduced in~\cref{sec:ez-mcts}.
While implementing basic \gls{mcts} is straightforward, some tweaks must be
introduced to guarantee a good throughput since \gls{mcts} is the main bottleneck of
training and self-play.
The core issue any \gls{mcts} implementation with imagined rollouts encounters
is the quick alternation between model inference on the GPU and using said inference
results for CPU work.
Node selection requires policy priors as well as state value and therefore
necessitates evaluating the policy and value heads $p$ and $v$.
Repeatedly running model inference on very small, unbatched inputs introduces major
kernel launch overhead and is therefore very inefficient.
The solution lies in batching as many GPU calls as possible.
This can be achieved for \gls{mcts} by stepping $n$ environments in parallel,
reducing the kernel launch overhead by a factor of $n$.
In practice, we implement this by keeping a pool of $n$ active environments and stepping
them in parallel while resetting any environments that have terminated.

Another bottleneck is host-device transfers between CPU and GPU.
For example, when evaluating multiple model heads, we concatenate their outputs
and transfer the result to the host in a single call instead of multiple calls.
Other general optimizations include reusing buffers as much as possible and precalculating
tables like the pUCT scaling term introduced in~\cref{eq:ez-puct}.
We also make use of rayon for simple task parallelism during backup.

\subsubsection{Replay Buffer}

The replay buffer stores past episodes including observations, actions and rewards,
as well as \gls{mcts} results like the root visit count distribution.
Because it must be efficiently sampleable for training, we use a struct-of-arrays representation
over ring buffers.
Since episodes may have different lengths, we additionally keep metadata such
as the offset of an episode into the data arrays as well as its length
and its initial observation.
The capacity of the replay buffer is fixed once at the beginning and can be
adjusted to fit the memory available on the system.

If the replay buffer is full, the oldest trajectory is evicted.
Sampling happens per step rather than per trajectory, in accordance with \gls{per}.
Freshly inserted episodes receive the largest priority observed so far, so that new data is
sampled at least once.
Priorities are updated by the reanalyze workers once the value error of a position is known.
Importantly, a sampled batch is assembled on the host and contains no GPU data.

The replay buffer is shared between the data workers that generate new
episodes, the rollout workers that sample raw batches, and the reanalyze
workers that re-run \gls{mcts} on the raw batches to produce finished batches.
It is guarded by a Mutex.

\subsubsection{Workers}

We make use of four different kinds of workers, of which a fixed number of instances
is spawned once in the beginning.
\Cref{fig:impl-ez-workers} shows the four kinds and the channels between them.
They communicate via message passing except for the replay buffer which is shared via a Mutex.
In general, each worker has access to a shared task queue on which it blocks and waits for
new tasks.
Once a task is finished, it publishes its result to a result queue that can then be consumed by other workers.
Since we assume the GPU to always be the bottleneck for training as well as data generation,
we run the workers in two phases: data collection and training.

\paragraph{Data workers}
They have access to a copy of a recent model and generate new episodes by interacting with the environment.
Each data worker contains a pool of active games that it steps in parallel using batched \gls{mcts}.
Once an episode is finished, it inserts the result into the replay buffer.
The learner generates a fixed number of data worker tasks so that the workers stop generating
new data after a limit is reached.

\paragraph{Rollout workers}
They sample the replay buffer, prepare raw batches and forward them to the reanalyze workers.
The number of raw batches is restricted by the learner in a similar manner as for the data workers.

\paragraph{Reanalyze workers}
They consume raw batches that have been prepared by the rollout workers and re-run \gls{mcts}
using the updated model parameters.

\paragraph{GPU workers}
All GPU work is done by GPU workers.
There are two reasons for this.
First, we found the Burn implementation to be extremely slow when relying on its internal
task queue due to contention.
Second, in its current implementation, Burn spawns a new CUDA stream for each thread.
We found that sending data between threads causes a lot of unwanted tensor data synchronization,
which in turn breaks kernel fusion, causing a major slowdown.
Wrapping all GPU work in an explicit GPU worker gives better control over queue
contention and data synchronization.
It also makes it possible to add a second device in the future by spawning a second worker.

%!TEX root = ../thesis.tex
% Thread and channel layout for \cref{sec:impl-efficientzero}: the four worker
% kinds, the one piece of shared mutable state they contend for, and the learner
% that drives them phase by phase. The algorithm-level counterpart is
% \cref{fig:ez-workers}, which draws the same loop as data flow; this one draws
% what is actually spawned and how the threads talk to each other.
%
% Layout rules:
%  * Three bands. The middle band is the pipeline in execution order, left to
%    right --- data workers, replay buffer, rollout workers, reanalyze workers
%    --- so every stage-to-stage edge is a short horizontal arrow between
%    neighbours. The GPU workers are the band above (every device request goes
%    straight up into it), the learner the band below (every task and status goes
%    straight down).
%  * Putting the replay buffer *inside* the row rather than into a band of its
%    own is what makes the figure crossing-free: it leaves exactly one
%    non-neighbour edge, the priority write-back from the reanalyze workers,
%    which is routed through the top lane at y=1.4 --- and the one column it
%    flies over, the rollout workers, is the one column with no GPU edge. Move
%    the buffer out of the row and that edge has nowhere left to go.
%  * Columns are 3.8cm apart and boxes are 2.6cm wide, so an edge label sitting
%    in a gap has 1.2cm (34pt) of room: at \scriptsize that is about nine
%    characters per line, which is why the in-row labels are stacked and terse.
%    "updated priorities" is single-line because it rides the long detour lane
%    rather than a gap. The picture is 14cm wide against a 14.56cm \textwidth
%    --- widen the columns before lengthening any in-row label.
%  * Blue = a worker kind, and there are deliberately exactly four of them.
%    Grey = shared state, white = the orchestrating thread.
%  * Solid arrows carry data or requests, dashed arrows carry parameters.
%  * The phase divider sits at x=5.7, which is exactly where the "sampled
%    positions" label already sits, and that label is wider than the 34pt gap
%    it nominally occupies. So the divider is two segments with the label in
%    between rather than one line behind a fill=white label: a fill wide enough
%    to hide the line also reaches into the two boxes either side, and one
%    anchored tight enough not to would eat the arrow it labels. The upper
%    segment starts at samplbl.north so it tracks the label if it is reworded.
%  * The divider is drawn after the pipeline arrows but before the priority
%    lane, so the lane crosses over it rather than under.
\begin{figure}[htbp]
  \centering
  \begin{tikzpicture}[
      >={Stealth[round]}, font=\small,
      worker/.style = {draw, rounded corners=2pt, fill=blue!8,
                       minimum height=15mm, minimum width=26mm, align=center},
      store/.style  = {draw, fill=black!8,
                       minimum height=15mm, minimum width=26mm, align=center},
      band/.style   = {draw, rounded corners=2pt, fill=blue!8, align=center},
      engine/.style = {draw, rounded corners=1.5pt, fill=white, align=center,
                       font=\scriptsize, minimum width=56mm, minimum height=11mm},
      orch/.style   = {draw, rounded corners=2pt, fill=white, align=center,
                       minimum width=140mm, minimum height=12mm},
      flow/.style   = {->, thick, black!70},
      param/.style  = {->, thick, black!70, densely dashed},
      lbl/.style    = {font=\scriptsize, inner sep=2.5pt, align=center},
    ]
    % --- the device band: the only threads that touch the GPU ----------------
    \node[band, minimum width=140mm, minimum height=20mm] (gpu) at (5.7,3.6) {};
    \node at (5.7,4.28) {GPU workers};
    \node[engine] (inf)   at (2.4,3.25) {model inference\\
                                         read-only model snapshot};
    \node[engine] (train) at (9.0,3.25) {model training\\
                                         autodiff, optimizer, gradient steps};
    \draw[param] (train.west) -- node[lbl,above] {$\theta$} (inf.east);

    % --- the pipeline, in execution order ------------------------------------
    \node[worker] (data) at (0,0)    {Data\\workers\\[1pt]
                                      \scriptsize self-play episodes};
    \node[store]  (buf)  at (3.8,0)  {Replay\\buffer\\[1pt]
                                      \scriptsize Mutex-guarded};
    \node[worker] (roll) at (7.6,0)  {Rollout\\workers\\[1pt]
                                      \scriptsize sample batches};
    \node[worker] (rean) at (11.4,0) {Reanalyze\\workers\\[1pt]
                                      \scriptsize refresh targets};

    \draw[flow] (data) -- node[lbl,above] {finished\\episodes}  (buf);
    \draw[flow] (buf)  -- node[lbl,above] (samplbl) {sampled\\positions} (roll);
    \draw[flow] (roll) -- node[lbl,above] {raw\\batches}        (rean);

    % --- the phase divider ---------------------------------------------------
    % Cuts the row between the replay buffer and the rollout workers, which is
    % where the two phases actually part: everything from the rollout workers
    % rightwards runs only in the training phase. It stops short of the GPU band
    % and the learner --- both are live in both phases, so extending it through
    % them would claim a split that does not exist.
    \draw[black!35, densely dashed] (5.7,-1.6) -- (5.7,0);
    \draw[black!35, densely dashed] (samplbl.north) -- (5.7,2.15);
    \node[lbl, black!55, anchor=south east] at (5.6,2.2) {data phase};
    \node[lbl, black!55, anchor=south west] at (5.8,2.2) {training phase};

    % --- device requests: straight up, one per column that needs one ---------
    \draw[flow] (data.north) --
      node[lbl,right,pos=0.62] {batched\\MCTS} (data.north |- gpu.south);
    \draw[flow] ([xshift=7mm]rean.north) --
      node[lbl,left,pos=0.7] {target \gls{mcts},\\gradient step}
      ([xshift=7mm]rean.north |- gpu.south);

    % --- the one non-neighbour edge, routed through the top lane -------------
    \draw[flow,rounded corners=3pt]
      ([xshift=-7mm]rean.north) -- (10.7,1.4)
      -- node[lbl,fill=white] {updated priorities} (3.8,1.4) -- (buf.north);

    % --- the orchestrator: tasks down, statuses up ---------------------------
    \node[orch] (learner) at (5.7,-2.35) {Learner\\[1pt]
      \scriptsize one iteration = a data phase, then a training phase};

    \draw[<->,thick,black!70] (data.south) --
      node[lbl,right] {tasks,\\statuses} (learner.north -| data.south);
    \draw[flow] (learner.north -| roll.south) --
      node[lbl,right] {tasks} (roll.south);
    \draw[flow] (rean.south) --
      node[lbl,right] {statuses} (learner.north -| rean.south);
  \end{tikzpicture}
  \caption[The worker threads of the EfficientZero implementation]{%
    The workers spawned by the learner and how they interact.
    Data workers receive data generation tasks and save episodes in the replay buffer.
    Rollout workers receive rollout tasks, sample raw batches from the replay buffer
    and forward them to reanalyze workers.
    Reanalyze workers refresh the raw batches using the latest model and then use
    backpropagation to calculate losses and step the model optimizer.
    The dashed line separates the data collection and training phases of a single iteration.
  }
  \label{fig:impl-ez-workers}
\end{figure}


\subsubsection{Hooks}\label{hooks}

To allow better diagnostics and instrumentation, the implementation
includes a way to add various kinds of hooks to different parts of the learning process.
This makes it possible to build a visualization dashboard or persistence layer
on top of the learner without needing to modify the learner itself.
\Cref{tab:impl-ez-hooks} contains a list of the hooks and where they are fired.

\begin{table}[htbp]
  \centering
  \caption[The instrumentation hooks of the learner]{%
    The instrumentation hooks of the learner, where they are fired and what
    context information the callback receives.
  }
  \label{tab:impl-ez-hooks}
  \footnotesize
  \begin{tabularx}{\textwidth}{@{}llX@{}}
    \toprule
    Hook & Fired by & Context \\
    \midrule
    \code{on\_env\_step}
      & data worker
      & action, reward and \code{done} of one environment step \\
    \code{on\_finished\_trajectory}
      & data worker
      & the finished episode, as inserted into the replay buffer \\
    \code{on\_finished\_data\_collection}
      & data worker
      & transitions produced by one data task \\
    \code{on\_finished\_mcts}
      & GPU worker
      & tree count, iterations per tree, root value of every tree \\
    \code{on\_search\_roots}
      & GPU worker
      & root priors next to the visit counts they were improved into,
        the only point at which the two coexist \\
    \code{on\_sampled\_batch}
      & rollout worker
      & the raw batch and the occupancy of the replay buffer \\
    \code{on\_finished\_reanalyze}
      & reanalyze worker
      & duration of one reanalyze task \\
    \code{on\_reanalyzed\_batch}
      & reanalyze worker
      & the training-ready batch, its scalar value targets, the root value
        predictions and the priorities written back \\
    \code{on\_finished\_training\_batch}
      & GPU worker
      & the four loss terms and the total loss of one gradient step \\
    \addlinespace
    \code{on\_finished\_data\_iteration}
      & learner
      & model snapshot after data collection, replay buffer \\
    \code{on\_finished\_training\_iteration}
      & learner
      & model snapshot after training, replay buffer \\
    \code{on\_finished\_learner\_iteration}
      & learner
      & iteration index, replay buffer length \\
    \bottomrule
  \end{tabularx}
\end{table}

Evaluation runs outside the learner and has three hooks of its own, fired per
step, per episode and once per run.
Because evaluation is single-threaded, they are generic over the concrete
\code{EvaluationEnvironment} and may therefore read target state that the
observation does not carry.
